% !TEX root = ../algebra_02.tex

\section{Ранг набора векторов. Столбцовый и строчный ранг матрицы}

\begin{Definition}{Ранг набора векторов}{}
	\emph{Рангом} набора векторов $v_1, \ldots, v_n$ называется размерность их линейной оболочки:
	\[
	\mathrm{rk}(v_1, \ldots, v_n) = \dim \mathrm{Lin}(v_1, \ldots, v_n).
	\]
\end{Definition}

Пусть $k$ — поле, $A \in k^{m \times n}$.

\begin{Definition}{Столбцовый и строковый ранги}{}
	\emph{Столбцовым рангом} матрицы $A$ называется ранг системы её столбцов:
	\[
	\mathrm{rk}_{\mathrm{col}}(A) = \mathrm{rk}(A[{:},1], \ldots, A[{:},n]).
	\]
	\emph{Строковым рангом} матрицы $A$ называется ранг системы её строк:
	\[
	\mathrm{rk}_{\mathrm{row}}(A) = \mathrm{rk}(A[1,{:}], \ldots, A[m,{:}]).
	\]
\end{Definition}

\begin{Statement}{Инвариантность рангов}{}
	Столбцовый и строковый ранги матрицы $A$ не меняются при элементарных преобразованиях строк матрицы $A$.
\end{Statement}

\begin{proof}
	Любое элементарное преобразование строк задаётся умножением слева на обратимую матрицу $U \in M_m(k)$: то есть $A' = UA$.
	
	\textbf{Инвариантность столбцового ранга.} Покажем, что $\mathrm{rk}_{\mathrm{col}}(A') = \mathrm{rk}_{\mathrm{col}}(A)$, то есть что из линейной зависимости столбцов $A'$ следует линейная зависимость столбцов $A$ и наоборот.
	Пусть $\lambda_1 A'[{:},i_1] + \cdots + \lambda_\ell A'[{:},i_\ell] = 0$.
	Так как $A'[{:},j] = UA[{:},j]$, получаем $U(\lambda_1 A[{:},i_1] + \cdots + \lambda_\ell A[{:},i_\ell]) = 0$.
	Поскольку $U$ обратима, $\lambda_1 A[{:},i_1] + \cdots + \lambda_\ell A[{:},i_\ell] = 0$. Обратное рассуждение симметрично (с использованием $U^{-1}$).
	Следовательно, системы столбцов линейно зависимы и независимы одновременно, откуда $\mathrm{rk}_{\mathrm{col}}(A') = \mathrm{rk}_{\mathrm{col}}(A)$.
	
	\textbf{Инвариантность строкового ранга.} Элементарные преобразования строк — это именно перестановки строк, умножение строки на ненулевой скаляр и прибавление одной строки к другой.
	Каждое из этих преобразований выражает новые строки как линейные комбинации старых, то есть не меняет линейную оболочку. Отсюда $\mathrm{rk}_{\mathrm{row}}(A') = \mathrm{rk}_{\mathrm{row}}(A)$.
\end{proof}

\begin{Remark}{}{}
	Если матрица $A$ элементарными преобразованиями строк приведена к ступенчатому виду с $r$ ненулевыми строками, то её строковый ранг равен $r$.
\end{Remark}
